\input{../../cm10192}
%\input{../../cm10192-handout}

\author[Lecture S.2.]{Lecture S.2.: Introduction to Linux/Unix Part 2
\\\begin{center}\input{../../copyright}\end{center}}


\AtBeginSection[]
{
  \begin{frame}<beamer>
    \frametitle{Outline}
    \tableofcontents[currentsection,currentsubsection,hideothersubsections]
  \end{frame}
}

%\AtBeginSubsection[]
%{
%  \begin{frame}<beamer>
%    \frametitle{Outline}
%    \tableofcontents[currentsection,currentsubsection,hide]
%  \end{frame}
%}

\logo{\vbox{\hbox to 1.1cm{\hfil\includegraphics[width=1cm,height=1.5cm]{../../computer}}}}



%------------------------------------------------------------------------

\begin{document}
\maketitle

\begin{frame}\frametitle{Outline}
\tableofcontents[pausesections,hidesubsections]
\end{frame}

\begin{frame}\frametitle{Resources}
\begin{itemize}
\item Unix for Beginners. Dirk Vermeir
\item \url{http://osl.iu.edu/~lums/swc/www/index.html}
\end{itemize}
\end{frame}

%\section{Questions \& Answers}

%\begin{frame}\frametitle{Q \& A}
%\begin{center}
%\includegraphics[height=7cm,keepaspectratio]{question}
%\end{center}
%\end{frame}

\section{Linux Part 2}
\begin{frame}{Previous Commands}
\rowcolors[]{1}{\lightrow}{\darkrow}
\begin{tabular}{|l|l|l|}
pwd & rm & cat\\
ls & chmod & whoami\\
cd & wc & passwd\\
mkdir & more & users\\
cp & less & groups\\
mv & diff & 
\end{tabular}
\end{frame}

\subsection{Redirecting IO}

\begin{frame}{Redirecting Input and Output}
\begin{itemize}
\item A running program is called a process
\begin{itemize}
\item By default, every process has three connections to the outside world:
\begin{itemize}
    \item Standard input (stdin): connected to the keyboard
    \item Standard output (stdout): connected to the screen
    \item Standard error (stderr): also connected to the screen (Used for error messages)
\end{itemize}
\end{itemize}
\end{itemize}
\end{frame}

\begin{frame}[fragile]\frametitle{Redirecting Input and Output II}
You can tell the shell to connect standard input and standard output to files instead
\begin{itemize}
    \item \lstinline!command < inputFile! reads from inputFile instead of from the keyboard
\begin{itemize}
          \item Don't need to use this very often, because most Unix commands let you specify the input file (or files) as command-line arguments
\end{itemize}
    \item \lstinline!command > outputFile! writes to outputFile instead of to the screen
\begin{itemize}
          \item Only “normal” output goes to the file, not error messages
\end{itemize}
    \item \lstinline!command < inputFile > outputFile! does both
\end{itemize}
\end{frame}

\begin{frame}[fragile]\frametitle{Redirecting Input and Output III}
\begin{itemize}
\item Example: save number of words in all text files to words.len
\codesmall
\codelist
\begin{lstlisting}[linewidth=5cm]
\$ wc -w *.txt > words.len
\end{lstlisting}
\item Nothing appears, because output is being sent to the file words.len
\reslist
\begin{lstlisting}[linewidth=5cm]
\$ ls -t
words.len       venus.txt       earth.txt
\$ cat words.len
 9  earth.txt
 9  venus.txt
 12 words.len
 30 total
\end{lstlisting}
\codenormal
Try typing \lstinline!cat > junk.txt!
\begin{itemize}
    \item No input file specified, so cat reads from the keyboard
    \item Output sent to a file
    \item Voila: the world's dumbest text editor
\end{itemize}
\end{itemize}
\end{frame}

\subsection{Pipes}

\begin{frame}[fragile]\frametitle{Pipes I}
\begin{itemize}
\item Suppose you want to use the output of one program as the input of another
\begin{itemize}
    \item E.g., use \lstinline!wc -w *.*! to count the words in some files, then \lstinline!sort -n! to sort numerically
\end{itemize}
\item Option 1: send output of first command to a temporary file, then read from that file
\codelist
\begin{lstlisting}[linewidth=5cm]
wc *.txt > temp
sort -n < temp
\end{lstlisting}
\end{itemize}
\end{frame}


\begin{frame}[fragile]\frametitle{Pipes II}
\begin{itemize}
\item Option 2: use a \keytext{pipe} to connect the two programs
\begin{itemize}
    \item Written as \lstinline!"|"!
    \item Tells the operating system to send what the first program writes to its stdout to the second program's stdin
\end{itemize}
\codelist
\begin{lstlisting}[linewidth=5cm]
wc -w *.* | sort -n
9 earth.txt
9 venus.txt
12 words.len
30 total
\end{lstlisting}
\item More convenient (and much less error prone) than temporary files
\end{itemize}
\end{frame}

\begin{frame}[fragile]\frametitle{Pipes III}
\begin{itemize}
\item Can chain any number of commands together
\begin{itemize}
    \item And combine with input and output redirection
    \item \codelist \codesmall 
\begin{lstlisting}[linewidth=7cm]
      wc *.txt | sort -n | head -5 > shortest.files
\end{lstlisting}
\end{itemize}
\item Any program that reads from standard input and writes to standard output can use redirection and pipes
\begin{itemize}
    \item Programs that do this are often called \keytext{filters}
    \item If you make your programs work like filters, you can easily combine them with others
    \item A combinatorial explosion of goodness
\end{itemize}
\end{itemize}
\end{frame}

\subsection{Environment Variables}

\begin{frame}[fragile]\frametitle{Environment Variables}
\begin{itemize}
\item Like any other program, the shell has variables
\item Since they define a user's environment, they are usually called environment variables
\item Usually all upper case
\item Type set at the command prompt to get a listing: (a cut down version)\\
\codelist
\begin{lstlisting}[linewidth=7cm]
\$ set
BASH=/bin/bash
HOME=/home/mdv
HOSTNAME=alis-tablet1
\end{lstlisting}
\end{itemize}
\end{frame}

\begin{frame}[fragile]\frametitle{Environment Variables II}
\begin{itemize}
\item Get a particular variable's value by putting a "\$" in front of its name
\begin{itemize}
\item E.g., the shell replaces "\$HOME" with the current user's home directory
\item Often use the echo command to print this out
\codelist
\begin{lstlisting}[linewidth=5cm]
$ echo $HOME
/home/mdv
\end{lstlisting}
\item Question: why must you type echo \$HOME, and not just \$HOME?
\end{itemize}
\item To set or reset a variable's value temporarily, use this:
\codelist
\begin{lstlisting}[linewidth=5cm]
\$ export VARNAME=value
\end{lstlisting}
\begin{itemize}
    \item Only affects the current shell (and programs run from it)
\end{itemize}
\end{itemize}
\end{frame}

\begin{frame}\frametitle{Environment Variables III}
\begin{itemize}
\item To set a variable's value automatically when you log in, set it in ~/.bashrc
\begin{itemize}
    \item Remember, "~" is a shortcut meaning “your home directory”
\end{itemize}
\end{itemize}
\end{frame}

\begin{frame}\frametitle{Environment Variables IV}
\rowcolors[]{1}{\lightrow}{\darkrow}
\footnotesize{
\begin{tabular}{|l|l|l}
\textbf{Name} & \textbf{Typical Value} & \textbf{Notes}\\
DISPLAY & 0:0 & \begin{minipage}{4cm}The display variable used for X11 graphics.\end{minipage}\\
HOME & /home/mdv &	The current user's home directory\\
HOMEDRIVE & C:	& \begin{minipage}{4cm}The current user's home drive (Windows only)\end{minipage}\\
HOSTNAME &	"alis-tablet"	 & This computer's name\\
HOSTTYPE &	"i486"	& What kind of computer this is\\
OSTYPE &	"linnux-gnu"	& What operating system is running\\
PATH	& \begin{minipage}{3cm}"/usr/local/sbin: \\/usr/local/bin: \\/usr/sbin:/usr/bin: \\/sbin:/bin:/usr/bin/X11: /usr/games"\end{minipage} & Where to look for programs\\
PWD	& /home/mdv/Desktop/ &	Present working directory \\
SHELL	& /bin/bash	& What shell is being run\\
TEMP	& /tmp &	Where to store temporary files\\
USER	& "mdv"	& The current user's ID
\end{tabular}}
\end{frame}

\begin{frame}\frametitle{How the Shell Finds Programs}
\begin{itemize}
\item The most important of these variables is PATH
\begin{itemize}
    \item The search path that tells the shell where to look for programs
    \item When you type a command like tabulate, the shell:
\begin{itemize}
          \item Splits \$PATH on colons to get a list of directories
          \item Looks for the program in each directory, in left-to-right order
          \item Runs the first one that it finds
\end{itemize}
\end{itemize}
\item Example
\begin{itemize}
    \item PATH is /home/mdv/bin:/usr/local/bin:/usr/bin:/bin:/Python24
    \item Both /usr/local/bin/tabulate and /home/mdv/bin/tabulate exist
    \item /home/mdv/bin/tabulate will be run when you type tabulate at the command prompt
    \item Can run the other one by specifying the path, instead of just the command name
\end{itemize}
\end{itemize}
\end{frame}

\begin{frame}[fragile]\frametitle{How the Shell Finds Program}
\begin{itemize}
\item Warning: it is common to include . in your path
\begin{itemize}
    \item This allows you to run a program in the current directory just by typing whatever, 
instead of ./whatever
    \item  But it also means you can never be quite sure what program a command will invoke
    \item  Though you can use the command \lstinline!which program_name!, which will tell you
\end{itemize}
\item Common entries in PATH include:
\begin{itemize}
    \item /bin, /usr/bin: core tools like ls (Note: the word “bin” comes from “binary”, which is geekspeak for “a compiled program”)
    \item  /usr/local/bin: optional (but common) tools, like the gcc
    \item \$HOME/bin: tools you have built for yourself
\end{itemize}          
\end{itemize}
\end{frame}

\subsection{Some More Tools}

\begin{frame}{Basic Tools I}
\rowcolors[]{1}{\lightrow}{\darkrow}
\footnotesize{
\begin{tabular}{|l|l}
man &	Documentation for commands. \\
cat &	Concatenate and display text files.\\
cd &	Change working directory.\\
chmod &	Change file and directory permissions.\\
clear &	Clear the screen.\\
cp &	Copy files and directories.\\
date &	Display the current date and time.\\
diff &	Show differences between two text files.\\
echo &	Print arguments.\\
env &	Show environment variables.\\
head &	Display the first few lines of a file.\\
ls &	List files and directories.\\
mkdir &	Make directories.\\
\end{tabular}}
\end{frame}

\begin{frame}{Basic Tools II}
\rowcolors[]{1}{\lightrow}{\darkrow}
\footnotesize{
\begin{tabular}{|l|l}
more &	Page through a text file.\\
mv &	Move (rename) files and directories.\\
od &	Display the bytes in a file.\\
passwd &	Change your password.\\
pwd & Print current working directory.\\
rm & Remove files.\\
rmdir &	Remove directories.\\
sort &	Sort lines.\\
tail &	Display the last few lines of a file.\\
uniq &	Remove duplicate lines.\\
wc &	Count lines, words, and characters in a file.\\
which &	locate a command\\
\end{tabular}}
\end{frame}

\begin{frame}[fragile]\frametitle{Compressing and Decompressing}
\begin{itemize}
\item If large files have to be stored, emailed or transported, it is easier to \keytext{compress} them, making the file size smaller. 
\item When needed later they can be restored or \keytext{decompressed}
\item A variety of compressing/decompressing algorithms exists.
\item A simple example: the sequence "aaaaa" could for example be represented more compactly as "5a"
\item The most common in linux systems are "gzip - gunzip","bzip2 - bunzip2" and "zip", which gives you the extensions ".gz",".bz2",".zip"
\end{itemize}
\reslist
\codesmall
\begin{lstlisting}
-rw-r--r-- 1 mdv mdv 514826 2007-10-11 00:01 lecture2.pdf
gzip lecture2.pdf
gunzip lecture2.pdf.gz
-rw-r--r-- 1 mdv mdv 382068 2007-10-11 00:01 lecture2.pdf.gz
bzip2 lecture2.pdf
-rw-r--r-- 1 mdv mdv 380635 2007-10-11 00:01 lecture2.pdf.bz2
bunzip2 lecture2.pdf.bz2
\end{lstlisting}
\end{frame}

\begin{frame}[fragile]\frametitle{Creating Archives}
\begin{itemize}
\item For storing purposes, emailing or transporting a number of files, it is sometimes more convenient to be able to bundle all these files together as one big file. Such a file is called an \keytext{archive}
\item In linux we use a program called "tar", so we call those archives tarfiles.
\item Format: \lstinline!tar -flags pathname [pathname]!
\item Common flags: 
\rowcolors[]{1}{\lightrow}{\darkrow}
\footnotesize{\begin{tabular}{l|l}
 c & create an archive\\
 x & extract an archive\\
t & display what is an archive\\
 v & verbose, display the files being dealt with\\
 f & use filename for archive\\
 z & filter through gzip/gzunip \\
j & filter through bzip2/bunzip2\\
\end{tabular}}
\end{itemize}
\end{frame}

\begin{frame}[fragile]\frametitle{Creating Archives: Some examples}
\begin{itemize}
\item \lstinline!tar -cvfz myfiles.tar.gz lectures/!: creates an archive called myfiles.tar which is also being compressed using gzip and will include the directory lectures and all files that are in the directory lectures.
\item \lstinline!tar -tvf myfiles.tar!: this command will list all files that are stored in the archive myfiles.tar
\item \lstinline!tar -xvfz myfiles.tar.gz!: this will command will decompress the file myfiles.tar.gz and will then extract the archive in the current directory. This will result in a directory lectures being created (if this directory does not exist) where all the files of archive are going to be put.
\end{itemize}
\end{frame}

\begin{frame}{More Advanced Tools II}
\rowcolors[]{1}{\lightrow}{\darkrow}
\footnotesize{
\begin{tabular}{|l|l}
du	& Print the disk space used by files and directories.\\
bzip    & Compress a file\\
bunzip  & Decompresses a file\\
find	& Find files that match a pattern.\\
grep	& Print lines matching a pattern.\\
gunzip	& Uncompress a file.\\
gzip	& Compress a file.\\
lpr	& Send a file to a printer.\\
lprm	& Remove a print job from a printer's queue.\\
lpq	& Check the status of a printer's queue.\\
ps	& Display running processes.\\
tar	& Archive files.\\
which	& Find the path to a program.\\
who	& See who is logged in.\\
xargs	& Execute a command for each line of input.
\end{tabular}}
\end{frame}

\begin{frame}{The Last Linux Slide: Making life Easier}
\begin{itemize}
\item When halfway typing a command on the command-line, try pressing the tab and see what happens.
\item After typing and executing a command on the command-line, try the arrow keys and see what happens
\end{itemize}
\end{frame}

\section{Problem Solving}

\begin{frame}[fragile]\frametitle{Scope}
\begin{multicols}{2}
\codelist
\begin{lstlisting}[linewidth=8cm]
x = 7

def scopetest(param):
   x = param
   print x
   if x == 0:
      return "yes"
   elif (x % 2) == 0:
      x = param / 2 - 1
      scopetest(x)
   else:
      x = x + 1
      scopetest(x)	

scopetest(x)	
\end{lstlisting}
\mbox{}\\
\mbox{}\\
\mbox{}\\
\mbox{}\\
\mbox{}\\
\onslide<2| handout:0| trans:0| article:0>
\reslist
\begin{lstlisting}[linewidth=2cm]
7
8
3
4
1
2
0
\end{lstlisting}
\end{multicols}
\end{frame}

\begin{frame}{How would you start on these?}
\begin{enumerate}[<+->]
\item Write a program that implements multiplication without using the use of \keytext{*}
\item Write a program that finds all prime dividers of a given number
\item Python has a module called random with a function random() providing random numbers between 0.0 and 1.0. How can you get numbers between 1 and 6 inclusive
\end{enumerate}
\end{frame}

\begin{frame}{And on these?}
\begin{enumerate}[<+->]
\item Given an ordered tree containing numbers, how would you go about finding a given number?
\item How would you go about writing a program that implements the game of Yahtzee?
\item Challenge! Global variables are not always a good idea. Everybody can access them. Suppose you want to limit access to a particular variables in such a way that can accessed in way you as a programmer intended.
\end{enumerate}
\end{frame}

\end{document}
